% Actividad 9 — Cuestionario evaluativo (standalone)
\documentclass[12pt]{article}
\usepackage[utf8]{inputenc}
\usepackage[T1]{fontenc}
\usepackage[spanish]{babel}
\usepackage{geometry}
\geometry{a4paper,margin=2.5cm}
\usepackage{enumitem}
\usepackage{parskip}
\usepackage{hyperref}
\title{Actividad 9 — Garantías Constitucionales: Cuestionario evaluativo}
\date{06 de septiembre de 2026}
\begin{document}
\maketitle
\section*{Instrucciones}
Fecha de entrega: 06 de septiembre de 2026, 23:55 CST. No se aceptan entregas extemporáneas.

Valor: 30 puntos (30/100 de la calificación final).

Alcance: tema 4 — Medios de control constitucional (4.1 a 4.6). El examen consta de 10 reactivos de opción múltiple o verdadero/falso. Cada reactivo vale 3 puntos.

\begin{itemize}
  \item Responda las 10 preguntas en este archivo y entregue el\pdf{} generado.
  \item Respuestas justificadas cuando se solicite (máx. 2 líneas por justificación).
  \item Cumplir con citas si corresponde; usar la bibliografía del curso (garantias-constitucionales.bib) según APA 7ª edición.
  \item No se permiten entregas extemporáneas.
\end{itemize}

\section*{Cuestionario (10 reactivos)}
\begin{enumerate}
  \item (V/F) El juicio de amparo procede únicamente contra actos de autoridad y no puede utilizarse para controlar la constitucionalidad de una ley en abstracto.
  \item (Múltiple) ¿Cuál de los siguientes mecanismos permite a la Suprema Corte declarar la invalidez general de una norma?
    \begin{enumerate}[label=\Alph*.]
      \item Juicio de amparo
      \item Controversia constitucional
      \item Recurso de revisión
      \item Acción de inconstitucionalidad
    \end{enumerate}
  \item (V/F) El control de convencionalidad exige que los jueces nacionales verifiquen la conformidad de las normas internas con los tratados internacionales de derechos humanos cuando resulten aplicables.
  \item (Múltiple) Según la jurisprudencia contemporánea de la SCJN, el parámetro de regularidad constitucional incluye:
    \begin{enumerate}[label=\Alph*.]
      \item Solo la Constitución
      \item Constitución y leyes generales
      \item Constitución, tratados de derechos humanos y su interpretación judicial
      \item Solo tratados internacionales
    \end{enumerate}
  \item (V/F) Cuando existe una restricción expresa en la Constitución que limita un derecho reconocido en un tratado, el juez debe aplicar siempre la norma constitucional sin considerar el principio pro persona.
  \item (Múltiple) El método de tres pasos para valorar jurisprudencia interamericana señala entre sus etapas:
    \begin{enumerate}[label=\Alph*.]
      \item Aplicar automáticamente la regla interamericana
      \item Verificar la identidad de razones y armonizar cuando sea posible
      \item Ignorar la jurisprudencia interamericana si México no fue parte
      \item Sustituir la Constitución por la interpretación interamericana
    \end{enumerate}
  \item (V/F) El control difuso habilita a cualquier juez para inaplicar una norma contraria a la Constitución en el caso concreto.
  \item (Múltiple) ¿Cuál es el referente legal primario para analizar la procedencia del juicio de amparo?
    \begin{enumerate}[label=\Alph*.]
      \item Ley de Amparo
      \item Código Civil
      \item Reglamento Universitario
      \item Ley Federal del Trabajo
    \end{enumerate}
  \item (V/F) La sentencia Radilla Pacheco y su seguimiento jurisdiccional tuvieron influencia en la exigencia de control de convencionalidad en los órganos judiciales mexicanos.
  \item (Múltiple) En caso de conflicto irreductible entre una interpretación interamericana más favorable y una restricción expresa constitucional, la práctica recomendada exige:
    \begin{enumerate}[label=\Alph*.]
      \item Aplicar siempre la interpretación interamericana
      \item Aplicar la restricción constitucional sin examen adicional
      \item Evaluar proporcionalidad y buscar la solución que más proteja al individuo, documentando la motivación
      \item Remitir el asunto al Congreso
    \end{enumerate}
\end{enumerate}

\section*{Referencias mínimas}
Consúltese como bases obligatorias: Constitución Política de los Estados Unidos Mexicanos; Ley de Amparo; Jurisprudencia de la SCJN y la bibliografía del curso (garantias-constitucionales.bib).

\vfill
\noindent\textit{Formato de entrega y citas: APA 7ª edición.}
\end{document}
